% !TeX root = ../main.tex
\section{Уравнения Максвелла}
\label{sec:maxwell}

Электромагнитное поле описывается
системой уравнений Максвелла.

В дифференциальной форме:

\begin{align}
\nabla \cdot \mathbf{D} &= \rho,\\
\nabla \cdot \mathbf{B} &= 0,\\
\nabla \times \mathbf{E}
    &= -\frac{\partial \mathbf{B}}{\partial t},\\
\nabla \times \mathbf{H}
    &= \mathbf{J}
    + \frac{\partial \mathbf{D}}{\partial t}.
\end{align}

Здесь $\mathbf{E}$ и $\mathbf{H}$ ---
напряжённости электрического
и магнитного полей.

Векторы $\mathbf{D}$ и $\mathbf{B}$
описывают электрическую индукцию
и магнитную индукцию соответственно.

В линейной изотропной среде:

\begin{align}
\mathbf{D} &= \varepsilon \mathbf{E},\\
\mathbf{B} &= \mu \mathbf{H}.
\end{align}

Для проводящей среды в простейшей
омической модели:

\begin{equation}
\mathbf{J}_{\mathrm{cond}}
=
\sigma \mathbf{E}.
\end{equation}

Подробное рассмотрение классической
электродинамики представлено
в работе~\cite{jackson1999}.
